\begin{center}
  \large\textbf{ABSTRAK}\\[2ex]
  \large\textbf{\tatitle{}}
\end{center}

\addcontentsline{toc}{chapter}{ABSTRAK}

\vspace{2ex}

\begingroup
% Menghilangkan padding
\setlength{\tabcolsep}{0pt}

\noindent
\begin{tabularx}{\textwidth}{l >{\centering}m{2em} X}
  Nama Mahasiswa / NRP   & : & \name{} / \nrp        \\

  Department & : & \department{} - Fakultas \faculty      \\

  Dosen Pembimbing        & : & 1. \advisor{}   \\
                    &   & 2. \coadvisor{} \\
\end{tabularx}
\endgroup

Pencatatan kehadiran mahasiswa secara manual di lingkungan universitas sering kali tidak efisien dan rentan terhadap kecurangan. Meskipun sistem absensi berbasis pengenalan wajah otomatis telah banyak diimplementasikan, metode konvensional berbasis \textit{Centralized Learning} menimbulkan kekhawatiran serius karena mewajibkan pengiriman citra wajah mentah mahasiswa ke server terpusat. Hal ini memicu risiko tinggi kebocoran privasi data biometrik yang sensitif serta pemborosan \textit{bandwidth} komunikasi jaringan. Untuk mengatasi tantangan tersebut, penelitian ini mengusulkan sebuah sistem presensi terdistribusi berbasis \textit{Edge Computing} yang mengintegrasikan metode \textit{Federated Learning} dengan arsitektur model \textit{MobileFaceNet} dan optimasi algoritma \textit{FedProx}. Seluruh komponen sistem dikemas menggunakan teknologi kontainerisasi \textit{Docker} untuk dideploy pada klaster perangkat keras \textit{edge} heterogen yang terdiri dari Jetson Nano dan Raspberry Pi 3 Model B. Pendekatan ini memfasilitasi pelatihan model secara lokal dan kolaboratif pada perangkat \textit{edge} tanpa mentransmisikan data biometrik mentah mahasiswa keluar dari masing-masing perangkat. Pengujian sistem mencakup evaluasi metrik pelatihan, efisiensi sumber daya jaringan, serta keandalan inferensi presensi pada kondisi nyata (\textit{real-life}) dan video simulasi. Hasil eksperimen menunjukkan model global \textit{Federated Learning} mencapai akurasi sebesar 98,16\% dengan \textit{validation loss} 0,3599, setara dengan metode terpusat yang mencatatkan akurasi 96,93\%. Selain itu, metode terdistribusi ini menghemat penggunaan \textit{bandwidth} secara signifikan dengan hanya mengirimkan data parameter sebesar 159,46 MB dibandingkan metode terpusat sebesar 1152,34 MB. Pada pengujian inferensi \textit{real-life}, sistem mencatatkan tingkat keberhasilan operasional 100\% pada jarak 0,5 hingga 3 meter di bawah pencahayaan cukup dan minim, dengan akurasi klasifikasi stabil di atas 90\% menggunakan ambang batas kemiripan 0,7. Latensi pemrosesan lokal pada Jetson Nano berlangsung sangat responsif ($\sim$700 ms), sedangkan Raspberry Pi membutuhkan $\sim$3100 ms akibat keterbatasan kapasitas RAM. Penelitian ini menyimpulkan bahwa integrasi \textit{Federated Learning} dan \textit{Edge Computing} terbukti efektif dalam menghadirkan sistem presensi yang aman, responsif, dan hemat \textit{bandwidth} jaringan.

% Ubah kata-kata berikut dengan kata kunci dari tugas akhir
\textbf{Kata Kunci: \keywords{}}
